\section{Proofs}
\label{app:proofs}

Theorem~\ref{thm:discounted_self_balancing} (Self-Balancing Under Discounting)
is proved in the main text (Section~\ref{sec:discounted}). This appendix
provides full proofs for the remaining results.

\subsection{Proof of Proposition~\ref{prop:preemptive} (Preemptive Restriction Criterion)}
\label{app:proof_preemptive}

\begin{proof}
The actor evaluates two action sequences from state $G_t$:

\emph{Sequence 1: Inaction.} The capability tuple remains intact. At each
future time step $t' > t$, there is probability $p(\Pathway \mid a_k, G_{t'},
\WorldModel)$ that the exercise pathway is executed, producing contraction
$-|\Delta \volL(\Pathway)|$. The expected future contraction is:
\[
\Vdisc(\text{inaction}) = -\sum_{t'=t+1}^{H} \gamma^{t'-t} \cdot
p(\Pathway \mid a_k, G_{t'}, \WorldModel) \cdot |\Delta \volL(\Pathway)|
\]
Under the simplifying assumption that the exercise probability is constant
($p$ does not change with time) and the pathway is a one-shot event (once
exercised, it does not repeat):
\[
\Vdisc(\text{inaction}) \geq -\Risk(d_1, \ldots, d_m; a_k, t)
\]
where $\Risk$ is the concentration risk
(Definition~\ref{def:concentration_risk}).

\emph{Sequence 2: Restrict.} The actor removes keystone capability $d_j$
from $a_k$ at time $t$, producing immediate contraction
$\Delta \volL(G_t \mid \text{restrict } d_j) < 0$ (by Lemma~2 of
\citet{lasser2026gfm}). All exercise pathways requiring $d_j$ are
eliminated; pathways not requiring $d_j$ persist under both sequences:
\[
\Vdisc(\text{restrict}) =
    -|\Delta \volL(G_t \mid \text{restrict } d_j)|
    + \text{(residual risk from non-}d_j\text{ pathways)}
\]

The residual risk term is identical under both sequences and cancels
in the comparison.  Restriction is value-positive
($\Vdisc(\text{restrict}) > \Vdisc(\text{inaction})$) when:
\[
-|\Delta \volL(G_t \mid \text{restrict } d_j)| >
-\Risk_j(d_1, \ldots, d_m; a_k, t)
\]
i.e., when $\Risk_j > |\Delta \volL|$, where $\Risk_j$ sums only over
pathways whose prerequisite set includes $d_j$.
\end{proof}

\subsection{Proof of Proposition~\ref{prop:structure_value} (Structural Discovery Value)}
\label{app:proof_structure_value}

\begin{proof}
The actor compares the expected discounted return of two strategies from
state $G_0$.

\emph{Strategy 1: Discover edge $d_1 \leq d_2$.} At $t = 0$, the actor
commits the subsumption edge, producing immediate contraction
$\Delta \volL(G_0 \mid \text{edge}) \leq 0$ (by Proposition~5(c) of
\citet{lasser2026poset}). For $t > 0$, the actor plans under the correct
world model $\WorldModel_{\text{true}}$:
\[
\Vdisc(\text{discover}) = \Delta \volL(G_0 \mid \text{edge}) +
\mathbb{E}_{\tau \sim \pi^*_{\text{true}}}\!\left[
\sum_{t=1}^{H} \gamma^t \Delta \volL(G_t)\right]
\]

\emph{Strategy 2: Avoid discovery.} At $t = 0$, no action is taken
($\Delta \volL = 0$). For $t > 0$, the actor plans under the incorrect
model $\WorldModel_{\text{false}}$:
\[
\Vdisc(\text{avoid}) = 0 +
\mathbb{E}_{\tau \sim \pi^*_{\text{false}}}\!\left[
\sum_{t=1}^{H} \gamma^t \Delta \volL(G_t)\right]
\]

Each expectation is taken over the trajectory distribution induced by its
respective policy. The difference is:
\begin{align*}
&\Vdisc(\text{discover}) - \Vdisc(\text{avoid}) \\
&\quad = \Delta \volL(G_0 \mid \text{edge})
    + I(d_1 \leq d_2)
\end{align*}
where $I(d_1 \leq d_2)$ is the information value of structure
(Definition~\ref{def:info_value}). Since
$\Delta \volL(G_0 \mid \text{edge}) \leq 0$, the difference is positive
when $I(d_1 \leq d_2) > |\Delta \volL(G_0 \mid \text{edge})|$.

The information value $I$ is non-negative in expectation: the policy
optimized under the correct model achieves at least as high an expected
return as the policy optimized under the incorrect model, because the
correct model produces unbiased $\volL$ estimates while the incorrect
model systematically overestimates the volume of capability sets with
shared prerequisites.
\end{proof}

\subsection{Proof of Proposition~\ref{prop:anti_monopolar} (Anti-Monopolar Property)}
\label{app:proof_anti_monopolar}

\begin{proof}
We prove each part in sequence.

\emph{Part (a): Individual growth-rate collapse.} In a fully dominant state
(Definition~\ref{def:dominant}), $\Gind{k} \subseteq G_K^{\mathrm{ind}}$
for every non-member $a_k$. Consider a non-member interaction at time $t$:
$a_k$ communicates a capability description $d$ or demonstrates $d$ through
the observation channel. Since $d \in \Gind{k} \subseteq G_K^{\mathrm{ind}}
\subseteq \C(t)$, the capability is already a node in the poset.
By Proposition~5(b) of \citet{lasser2026poset}, adding a capability already
in $\C$ does not increase $|\C(t)|$. Therefore no non-member's individual
capabilities generate novel poset nodes, and the individual-capability
component of $r_{\mathrm{ext}}$ is zero. Novel cooperative capabilities
from cross-coalition interaction may persist
(Section~\ref{sec:substitution_coop}), making this a conservative bound.

\emph{Part (b): Trajectory comparison.} The actor at $t = 0$ compares two
strategies, each evaluated by $\Vdisc = \sum_t \gamma^t \Delta\volL(G_t)$
(Definition~\ref{def:discounted}). Under the linearized model, each
strategy produces a constant per-step $\Delta\volL$. We compare the
discounted sums $\Vdisc$ for $\gamma \in (0, 1)$.

\emph{Strategy D (dominate):} At $t = 0$, subsume remaining non-member
capabilities, gaining $\Delta\volL = \Delta_0 \geq 0$. For $t \geq 1$,
grow at internal rate: $\Delta\volL = r_K$ per step.
\[
\Vdisc^D = \Delta_0 + \frac{r_K \cdot \gamma}{1 - \gamma}
\]
\emph{Strategy Div (maintain diversity):} No subsumption bonus. For
$t \geq 0$, grow at full rate: $\Delta\volL = r_K + r_{\mathrm{ext}}$
per step.
\[
\Vdisc^{\mathrm{div}} = \frac{r_K + r_{\mathrm{ext}}}{1 - \gamma}
\]
The difference:
\begin{align*}
\Vdisc^{\mathrm{div}} - \Vdisc^D
  &= \frac{r_K + r_{\mathrm{ext}}}{1 - \gamma}
     - \Delta_0 - \frac{r_K \gamma}{1 - \gamma} \\
  &= r_K + \frac{r_{\mathrm{ext}}}{1 - \gamma} - \Delta_0
\end{align*}
This is positive when $r_{\mathrm{ext}} / (1 - \gamma) > \Delta_0 - r_K$.
\emph{Case 1:} If $\Delta_0 \leq r_K$, the left side is positive for all
$\gamma > 0$. \emph{Case 2:} If $\Delta_0 > r_K$, rearranging:
$1 - \gamma < r_{\mathrm{ext}} / (\Delta_0 - r_K)$, i.e.,
\[
\gamma > \gamma^* = 1 - \frac{r_{\mathrm{ext}}}{\Delta_0 - r_K}
    = \frac{\Delta_0 - r_K - r_{\mathrm{ext}}}{\Delta_0 - r_K}
\]
When $\Delta_0 - r_K \leq r_{\mathrm{ext}}$, $\gamma^* \leq 0$ and the
condition holds for all $\gamma > 0$. When $\Delta_0 - r_K >
r_{\mathrm{ext}}$, $\gamma^* \in (0, 1)$ and a sufficiently long planning
horizon makes diversity strictly preferable.

\end{proof}

Corollary~\ref{cor:a1_from_oi} (Static Protection from Observational
Individuation) and Corollary~\ref{cor:gradual} (Gradual Domination) are
proved in the main text (Sections~\ref{sec:obs_individuation}
and~\ref{sec:gradual} respectively), as their proofs are short and
self-contained.
